%==============================================================================
%  IEEE T-BIOM regular paper (D53): the threshold-transfer audit, two registered rounds.
%  Round 1 = the registered audit as revised for the simulated ICME panel (DFV_Conference/revision, audited numbers).
%  Round 2 = PREREGISTRATION_ROUND2.md (tag round2-freeze). Every round-2 number comes from a generated table
%  (analysis/round2/r2_paper_tables.py -> gen/r2_*.tex) or carries a "% src:" comment; \todo{} marks what is pending.
%  At most 10 pages (T-BIOM regular paper); supplement in supplement.tex. Build: pdflatex/bibtex/pdflatex x2.
%==============================================================================
\documentclass[journal]{IEEEtran}
\usepackage{cite}
\usepackage{amsmath,amssymb}
\usepackage{graphicx}
\graphicspath{{figures/}}
\usepackage{booktabs}
\usepackage{xcolor}
\usepackage{url}
\usepackage{microtype}

\newif\ifblind
\blindtrue   % T-BIOM allows double-anonymous submission at the authors' choice ([U] decision)
\newcommand{\todo}[1]{\textcolor{red}{[#1]}}
\newcommand{\gentab}[1]{\IfFileExists{gen/#1.tex}{\input{gen/#1}}{\todo{generated table \detokenize{#1} pending}}}
\ifblind
  \newcommand{\repo}{\url{https://anonymous.4open.science/r/ttg-audit}}
\else
  \newcommand{\repo}{\todo{public repository URL}}
\fi

\newcommand{\ttg}{\mathrm{TTG}}
\newcommand{\hter}{\mathrm{HTER}}
\newcommand{\frr}{\mathrm{FRR}}
\newcommand{\far}{\mathrm{FAR}}
\newcommand{\eer}{\mathrm{EER}}
\newcommand{\tauE}{\tau^{*}_{\mathrm{EER}}}
\newcommand{\tauF}[1]{\tau^{*}_{\mathrm{FRR}#1}}
\newcommand{\tauA}[1]{\tau^{*}_{\mathrm{FAR}#1}}

\begin{document}

\title{Thresholded Once, Deployed Elsewhere:\\ A Two-Round Pre-Registered Audit of Threshold Transfer in Speech Deepfake Detection}

\ifblind
\author{Anonymous T-BIOM Submission}
\else
\author{Hoan~Hoang\thanks{\todo{affiliation, e-mail}}}
\fi

\maketitle

\begin{abstract}
\todo{Abstract after the round-2 results (<= 250 words). Must state: the protocol card; round 1 (five LA19-trained
checkpoints; H2 holds for clips but not for speakers; H5 holds; a dev threshold costs AASIST/AASIST-L 4.2/4.6 HTER points on
LA19 eval; the error direction follows the policy; tie intervals); round 2 (three current detectors incl. ASVspoof 5-trained
and calibration-oriented; seed variance; silence-trimmed training; eight channel conditions; three-zone policy) with the
registered outcomes R2-H1..H4; the D50 sentence on H3/H6 (not supported; reported separately).}
\end{abstract}

\begin{IEEEkeywords}
speech deepfake detection, presentation attack detection, anti-spoofing, decision threshold, calibration, domain shift,
pre-registration, ISO/IEC 30107-3
\end{IEEEkeywords}

% ============================ I. INTRODUCTION ================================
\section{Introduction}
\label{sec:intro}
\IEEEPARstart{S}{peech} deepfake countermeasures (CMs) are ranked by the equal error rate (EER) and the minimum tandem
detection cost \cite{kinnunen18_tdcf}, both evaluated at the threshold that suits the evaluation data best
\cite{wang20_asvspoof2019,liu23_asvspoof2021}. A deployed CM needs more: a call-screening filter, a moderation queue or a
browser plug-in \cite{pascu25_edge_plugin} commits to a threshold in advance and decides on traffic it does not control.
The EER does not say how that threshold behaves there, and the difference can be large: Zhou and Wang
\cite{zhou26_eer_hides} report a detector with 0.21\% source EER that rejects 78.7\% of genuine In-the-Wild speech once its
source threshold is transferred.

Fixing a threshold on development data and reporting the half total error rate (HTER) on test data is established practice
in presentation-attack detection (PAD) \cite{chingovska12_lbp,korshunov16_crossdb} and biometric authentication
\cite{bengio04_hter}, and the loss it exposes is the calibration loss of application-independent evaluation
\cite{brummer06_calibration,vanleeuwen07_intro}. Speech CMs lack an audit built on that lineage: several models and
threshold policies, an error decomposition, speaker-aware inference, hypotheses fixed in advance, and a test of whether the
findings outlive the checkpoints they were found on.

This paper is that audit, in two registered rounds. Round~1 audits five released checkpoints trained on ASVspoof~2019 LA
(LA19). Round~2, registered before any of its scores existed, asks whether the round-1 findings hold for current detectors
(one trained on ASVspoof~5, one calibration-oriented, one 2024--25 self-supervised model), for retrained copies of one
round-1 recipe, for the same recipe trained without the leading-silence cue, and under eight more channel conditions.
We contribute:
\begin{itemize}
\item a protocol card for threshold-transfer audits (Table~\ref{tab:card}): the threshold transfer gap (TTG) as the
calibration loss of a development threshold, split from the discrimination loss that no threshold removes; five threshold
policies with their ISO/IEC 30107-3 reading; the interval of thresholds that tie on development data; speaker-level
inference; and a cost model with its break-even window;
\item a two-round pre-registered audit: round-1 confirmatory tests (H2, H5) and exploratory findings, and round-2
confirmatory replications of those findings on fresh models, runs and channels (R2-H1 to R2-H4);
\item a causal test of the silence explanation, by training the same recipe on silence-trimmed audio;
\item a three-zone (accept / review / reject) policy designed on development data alone, with what it sends to review
after transfer;
\item the harness, frozen thresholds, score files and analysis scripts (\repo).
\end{itemize}
Scores are fp32 on a desktop; deployment on devices is reported separately.

% ============================ II. RELATED WORK ================================
\section{Related Work}
\label{sec:related}
\emph{Evaluation practice.} ASVspoof ranks CMs by EER and minimum t-DCF \cite{wang20_asvspoof2019,liu23_asvspoof2021,
kinnunen18_tdcf}. ASVspoof~5 also reports actual costs at a fixed threshold, and its strongest systems pair low minimum
costs with actual costs near 1 \cite{wang24_asvspoof5}, a calibration failure that EER cannot show. The gap between actual
and minimum cost is the calibration loss of speaker recognition \cite{brummer06_calibration,vanleeuwen07_intro}. In PAD,
HTER at a development threshold is the standard protocol for face \cite{chingovska12_lbp} and voice
\cite{korshunov16_crossdb}, including across databases, where HTER reaches 50\% \cite{korshunov16_crossdb}; ISO/IEC 30107-3
names the two error rates APCER and BPCER \cite{iso30107_3}.

\emph{Generalization and artefacts.} In-the-Wild (ITW) \cite{muller22_in_the_wild} showed that CMs trained on ASVspoof data
degrade sharply on public figures' speech from online media. M\"uller et al.\ \cite{muller21_silence} found that the
duration of leading silence separates LA19 bona fide from spoofed speech, and that trimming it before training raised EER
from 3.6\% to 15.5\%. ASVspoof~5 reduces such shortcuts by post-processing (peak normalisation, random trimming of leading
and trailing non-speech) in its development and evaluation sets, but not in its training set \cite{wang25_asvspoof5_db}.
Calibration-oriented detectors such as frozen self-supervised features with logistic regression \cite{pascu24_calibrated}
report low expected calibration error; whether that survives a transferred threshold is one of our questions. Seed
variance of the same CM can be large \cite{wang21_comparative}.

\emph{Threshold transfer and selective decisions.} The closest precedent \cite{zhou26_eer_hides} freezes one SSL detector,
sets its threshold on the labelled LA19 evaluation set and transfers it to ITW and a DF21 subset. Our thresholds come from
development data, so no evaluation label informs them. Abstaining on uncertain inputs is the reject option of selective
classification \cite{geifman17_selective}; a three-zone policy is its operational form when the abstentions go to human
review.

\emph{Channels and codecs.} DF21 builds media compression into its evaluation \cite{liu23_asvspoof2021}. Neural codecs are
also a synthesis route \cite{wu24_codecfake}, so the label of codec-resynthesised bona fide speech is contested
\cite{xiao26_resynth}; ASVspoof~5 applies EnCodec to subsets of both classes as a channel \cite{wang25_asvspoof5_db}.

% ============================ III. PROTOCOL =====================================
\section{Protocol}
\label{sec:protocol}

\begin{table}[!t]
\caption{Protocol card for a threshold-transfer audit.}
\label{tab:card}
\centering
\footnotesize
\setlength{\tabcolsep}{3pt}
\begin{tabular}{@{}p{0.21\columnwidth}p{0.73\columnwidth}@{}}
\toprule
\emph{Inputs} & A CM that outputs a score; labelled source development data; target sets with speaker labels; one or more cost
models. \\
\emph{Frozen before any target score} & Threshold policies computed from development scores alone (EER; FRR budgets 1\% and
5\%; FAR budgets 1\% and 5\%), the estimator and its tie rule, the hypotheses and the resampling unit (speaker). \\
\emph{Report at each policy and target} & FRR (BPCER) and FAR (APCER) at the shipped threshold; HTER split into oracle HTER
and TTG, with speaker-level intervals; FRR and FAR at both ends of the interval of thresholds that tie on dev data; cost over
a sweep of priors, with its break-even window; the data and policy that set the threshold. \\
\emph{For triage} & The review share and the automatic-decision errors of a three-zone policy designed on dev data. \\
\bottomrule
\end{tabular}
\end{table}

\subsection{Scores, thresholds and measures}
Each CM maps a clip \(x\) to a score \(s(x)\), higher meaning more bona fide, and declares a spoof if \(s(x)<\tau\). Every
threshold is fixed on the source development set of the model's training data. \(\tauE\) is the crossing of dev FRR and
FAR, estimated by linear interpolation between adjacent unique dev scores. \(\tauF{1}\) and \(\tauF{5}\) are the strictest
thresholds whose dev FRR stays within 1\% and 5\%, and \(\tauA{1}\) and \(\tauA{5}\) the most lenient thresholds whose dev
FAR stays within 1\% and 5\% (BPCER at a fixed APCER in ISO/IEC 30107-3 terms). On a target, \(\frr(\tau)\) is the share of
bona fide clips declared spoof (BPCER), \(\far(\tau)\) the share of spoofs accepted (APCER), and
\(\hter(\tau)=\tfrac12[\frr(\tau)+\far(\tau)]\) \cite{bengio04_hter}. For a transferred threshold \(\tau^{*}\),
\begin{equation}
\ttg(\tau^{*}) \;=\; \hter(\tau^{*}) \;-\; \min_{\tau}\hter(\tau)
\label{eq:ttg}
\end{equation}
is the calibration loss of \(\tau^{*}\) at the equal-cost point \cite{brummer06_calibration,vanleeuwen07_intro}. The
minimum, the oracle HTER, is the part of the deployed error that no fixed threshold removes (discrimination), and TTG the
part that a better threshold would remove (transfer). We report both, and the transfer share \(\ttg/\hter(\tau^{*})\).
Because a near-separable development set pins \(\tauE\) with few trials, we also report the \emph{dev-equivalent
interval}: the connected run of thresholds around \(\tauE\) whose dev HTER exceeds \(\hter(\tauE)\) by at most one bona
fide trial (\(0.5/n_{\mathrm{bona}}\)), and the target FRR at its two ends.

\emph{Cost.} With \(p\) the bona fide prior weighted by costs, \(C(p,\tau)=[p\,\frr(\tau)+(1-p)\,\far(\tau)]/\min(p,1-p)\)
is 1 for the better trivial system, so a detector beats both trivial systems only for priors in
\(\far/(\far+1-\frr)<p<(1-\far)/(1-\far+\frr)\), its break-even window.

\emph{Three-zone policy.} A review band \([\tau_{\mathrm{lo}},\tau_{\mathrm{hi}})\) holding 10\% of the pooled dev scores,
centred in rank on \(\tauE\), sends scores inside it to review, rejects below it and accepts above it. On each target we
report the review share and the automatic errors, and compare with the same rule fitted to the labelled target and with a
label-free variant whose edges are re-anchored on the target's unlabelled scores.

% ============================ IV. STUDY DESIGN =====================================
\section{Study Design}
\label{sec:design}
\subsection{Round 1: released LA19-trained checkpoints}
% src: the ICME revision (DFV_Conference/revision/paper/main.tex Sec. III-B/C; audited by audit_claims.py).
Round~1 audits five published CMs: four primary models (AASIST and AASIST-L \cite{jung22_aasist}, RawNet2
\cite{tak21_rawnet2}, Res-TSSDNet \cite{hua21_res_tssdnet}) and a self-supervised reference, XLS-R+AASIST
\cite{tak22_xlsr_aasist} (Table~\ref{tab:roster}). Checkpoints are used as released; we vendor each author's code and keep
its score function, class order and input policy. A primary model must reproduce its published LA19 evaluation EER within
15\% (relative); the released RawNet2 is the ASVspoof~2021 DF baseline and is gated on DF21 against the organisers'
released scores (per-trial agreement \(3.7{\times}10^{-5}\)); XLS-R+AASIST has no published LA19 figure and enters no
hypothesis count. Every threshold is fixed on the LA19 development set (24,844 trials, 2,548 bona fide).

\begin{table}[!t]
\caption{Round-1 roster and development thresholds. EER on LA19 eval (\%): ours / published. Dev EER (\%) and \(\tauE\) on
LA19 dev. Last column: dev bona fide trials below, and dev spoof trials at or above, \(\tauE\).}
\label{tab:roster}
\centering
\footnotesize
\setlength{\tabcolsep}{3pt}
\begin{tabular}{@{}lrcrrc@{}}
\toprule
Model & Params & Eval EER & Dev EER & $\tauE$ & Dev trials \\
\midrule
AASIST & 298K & 0.83 / 0.83 & 0.510 & $-$1.09 & 13 / 113 \\
AASIST-L & 85K & 0.99 / 0.99 & 0.510 & $-$0.94 & 13 / 113 \\
RawNet2 & 17.6M & 4.60 / 5.64\textsuperscript{a} & 0.549 & $-$7.07 & 14 / 122 \\
Res-TSSDNet & 349K & 1.62 / 1.66 & 0.722 & $-$8.24 & 19 / 161 \\
XLS-R (ref.) & 315M & 0.15 / ---\textsuperscript{b} & 0.009 & $-$1.75 & 1 / 2 \\
\bottomrule
\end{tabular}
\\[2pt]
\parbox{\columnwidth}{\footnotesize\textsuperscript{a}Different training run, gated on DF21 instead.
\textsuperscript{b}None published, so ungated and outside every hypothesis count.}
\end{table}

The targets are the LA19 evaluation set (71,237 trials; 11 of 13 attacks unseen), In-the-Wild (31,779 clips of public
figures with 54 distinct speaker labels in its metadata), a verified subset of ASVspoof~2021 DF (455,508 trials; libFLAC~1.4
rejects 43\% of the public evaluation files and for 78,420 of them the only working decode fails the file's own MD5, so we
exclude those), and an EnCodec stress set: 1,000 LA19-eval clips (500 per class, all 67 speakers) resynthesised at 6\,kbps
\cite{defossez23_encodec}, whose bona fide clips stay bona fide (a genuine user's speech through a codec is still
genuine) \cite{wang25_asvspoof5_db,xiao26_resynth}. The registered hypotheses were H2 (on ITW at \(\tauE\), FRR exceeds FAR
for at least three of the four primary models) and H5 (EnCodec raises FRR more than FAR, relative to the clean copy, for
at least three), tested by one-sided class-stratified bootstraps (B = 1,000) with Holm correction; because clips are nested
in speakers, we added speaker-cluster bootstraps (B = 10,000) as a logged deviation \cite{field07_boot,bolle04_subsets}.
The hypotheses were registered on 19~July 2026, amended in August and September, and frozen on 28~September 2026 before any
ITW clip was scored; the frozen registration and the deviation log are in the repository. Two int8 hypotheses were also
registered; neither was supported (H3 for lack of usable int8 models, H6 on a clause the quantizer could never meet), and a
separate study reports both.

\subsection{Round 2: registered extension}
Round~2 was registered, with its analysis code, before any of its scores existed (tag \texttt{round2-freeze}; deviation log
entry DEV-9 with three amendments, each made before the work it governs ran). It adds three detectors outside every round-1
count: the ASVspoof~5 baseline AASIST (B02), trained on the un-post-processed ASVspoof~5 training set and thresholded on its
post-processed development set \cite{wang25_asvspoof5_db}; the calibration-oriented detector of Pascu et
al.\ \cite{pascu24_calibrated} (frozen XLS-R~2B \cite{babu21_xlsr} features, logistic regression), refit on LA19 train only so that LA19 dev
stays held out; and XLSR-Mamba \cite{xiao25_xlsrmamba}, an LA19-trained SSL model. It trains the round-1
AASIST-L recipe five times (seeds 1--5) as released and five times on audio whose leading and trailing non-speech is
removed by an energy VAD, keeping the recipe's final weight-averaged model. It adds eight channel conditions on the 1,000
H5 clips: Opus at 6 and 24 kbps, AMR-NB at 4.75 kbps, MP3 at 32 kbps, MUSAN noise \cite{snyder15_musan} at 20, 10 and 5 dB SNR,
and a mock-up phone impulse response \cite{jeub09_air,ko17_rirs}. ITW, DF21, MLAAD-5 \cite{mlaad24} and the
ASVspoof~5 evaluation set (for the ASVspoof~5 model) are targets. The confirmatory hypotheses are:
\begin{itemize}
\item \emph{R2-H1}, the policy sets the direction: on ITW, FRR > FAR at \(\tauF{1}\) and FAR > FRR at \(\tauA{5}\), for at
least two of the three new detectors;
\item \emph{R2-H2}, the LA19-eval transfer gap survives retraining: TTG at \(\tauE\) > 1\,pp for at least four of five
AASIST-L runs;
\item \emph{R2-H3}, removing the silence cue from training lowers ITW false rejection (and, separately, ITW EER), with a
higher LA19-dev EER as the manipulation check;
\item \emph{R2-H4}, H5 generalised: each channel condition raises FRR more than FAR for at least three of the four primary
round-1 models.
\end{itemize}
All use speaker-cluster bootstraps (B = 10,000) or, for R2-H3, exact permutation tests over training runs, with Holm
correction \cite{holm79_holm}. A reading rule fixed in advance classifies each new detector as discrimination-limited,
transfer-limited, well-transferred or intermediate.

% ============================ V. RESULTS: ROUND 1 ======================================
\section{Round 1: Released Checkpoints}
\label{sec:round1}
% src: the ICME revision Sec. IV (results/b_thresholds.txt, a_cluster.txt, f_vad_*.txt, a5_cheapcheck.txt, e_cost.txt).
\begin{table*}[!t]
\caption{Round 1: operating behaviour at the frozen \(\tauE\) (fp32, \%). HTER = oracle HTER (discrimination) + TTG
(transfer, pp); TTG/HTER is the transfer share (\%). The CI of TTG is a speaker-cluster bootstrap (B = 1,000).}
\label{tab:op}
\centering
\footnotesize
\setlength{\tabcolsep}{5pt}
\begin{tabular}{@{}lrrrrrrlr@{}}
\toprule
Model & EER & FRR & FAR & HTER & Oracle & TTG & 95\% CI (speaker bootstrap) & TTG/HTER \\
\midrule
\multicolumn{9}{@{}l}{\emph{LA19 eval: seen corpus, unseen attacks} (67 speakers)} \\
AASIST & 0.83 & 0.16 & 9.83 & 5.00 & 0.82 & 4.18 & [3.74, 4.65] & 84 \\
AASIST-L & 0.99 & 0.10 & 11.14 & 5.62 & 0.97 & 4.64 & [3.99, 5.36] & 83 \\
RawNet2 & 4.60 & 0.50 & 11.48 & 5.99 & 4.25 & 1.74 & [1.40, 2.24] & 29 \\
Res-TSSDNet & 1.62 & 0.26 & 6.08 & 3.17 & 1.59 & 1.58 & [1.30, 1.89] & 50 \\
XLS-R (ref.) & 0.15 & 0.05 & 0.49 & 0.27 & 0.14 & 0.13 & [0.08, 0.18] & 46 \\
\midrule
\multicolumn{9}{@{}l}{\emph{In-the-Wild} (54 speaker labels)} \\
AASIST & 43.0 & 43.8 & 41.9 & 42.8 & 39.9 & 2.94 & [0.16, 7.25] & 7 \\
AASIST-L & 44.5 & 50.6 & 36.5 & 43.6 & 42.0 & 1.55 & [0.19, 5.39] & 4 \\
RawNet2 & 49.0 & 33.8 & 69.1 & 51.4 & 46.0 & 5.44 & [2.56, 9.11] & 11 \\
Res-TSSDNet & 34.2 & 50.3 & 19.0 & 34.6 & 33.9 & 0.73 & [0.07, 5.17] & 2 \\
XLS-R (ref.) & 10.5 & 49.9 & 0.4 & 25.2 & 10.2 & 14.96 & [11.42, 19.66] & 59 \\
\midrule
\multicolumn{9}{@{}l}{\emph{DF21, verified subset} (93 speakers)} \\
AASIST & 15.0 & 13.8 & 16.2 & 15.0 & 14.8 & 0.29 & [0.08, 1.92] & 2 \\
AASIST-L & 14.9 & 16.5 & 13.3 & 14.9 & 14.8 & 0.03 & [0.02, 0.91] & 0 \\
RawNet2 & 19.7 & 33.7 & 5.6 & 19.6 & 19.1 & 0.58 & [0.04, 2.53] & 3 \\
Res-TSSDNet & 28.1 & 52.8 & 5.2 & 29.0 & 27.2 & 1.82 & [0.94, 3.61] & 6 \\
XLS-R (ref.) & 2.0 & 14.6 & 0.2 & 7.4 & 2.0 & 5.40 & [3.23, 7.77] & 73 \\
\bottomrule
\end{tabular}

\end{table*}

\emph{Unseen attacks.} On LA19 eval the genuine speech comes from the source corpus and only the attacks are new, so the
error that grows is false acceptance: the primary models accept 6.1--11.5\% of spoofs and reject 0.1--0.5\% of genuine
clips. AASIST and AASIST-L lose 4.18 and 4.64 HTER points to transfer (speaker-level 95\% CIs [3.74, 4.65] and
[3.99, 5.36]), 84\% and 83\% of their deployed HTER, which the EER hides (0.83\% against 5.00\% for AASIST). By HTER at the
development threshold, Res-TSSDNet (3.17\%) overtakes both, which beat it by EER.

\emph{Shifted targets.} On ITW and DF21 the genuine speech has shifted and FRR exceeds FAR for four of five models on each
target. For the primary models on ITW the dominant loss is discrimination (EER 34--49\%; transfer only 2--11\% of the
deployed HTER). The reference has the largest gap from its own optimum (59\% and 73\% of its HTER on ITW and DF21) yet the
fewest errors (HTER 25.2\% on ITW against 34.6--51.4\%).

\emph{The registered tests.} H2 is supported at the clip level (three of four) but not at the speaker level: resampling
ITW's 54 speaker labels, only Res-TSSDNet's FRR $-$ FAR excludes zero (+31.3\,pp [+14.3, +51.8]), and deleting a single
speaker moves AASIST's difference between $-0.5$ and +5.4\,pp; we do not claim H2. H5 is supported under both units:
EnCodec raises FRR and lowers FAR for every primary model (Res-TSSDNet rejects 68.4\% of resynthesised bona fide clips).

\emph{The policy sets the direction.} A threshold whose dev FRR is 1\% rejects 42.8--81.8\% of genuine ITW speech, while a
dev FAR of 5\% becomes 14.2--91.9\% FAR with FAR above FRR for all five models: the dominant error follows the policy, and
the registered grid, all on the strict side, tilted the registered finding toward false rejection.

\emph{Determinacy.} Dev is nearly separable, so \(\tauE\) is pinned by few trials. For the reference, every threshold in
\([-3.38, +2.14]\) ties with \(\tauE\) on dev within one bona fide trial, and its ITW FRR runs from 25.1\% to 73.0\% over
that interval.

\emph{Silence.} LA19 bona fide clips begin with a median 1.00\,s of silence against 0.38\,s for spoofs; ITW genuine clips
begin with 0.09\,s. Trimming both ends at test time raises the LA19-eval EER of all five models 6--34 fold, so each relies on
the cue; re-setting thresholds on trimmed dev lowers ITW FRR (9--38\,pp) but raises FAR (3--34\,pp) and moves ITW EER by
only \(-4.0\) to \(+3.1\)\,pp. Whether the cue in \emph{training} drives the ITW behaviour is R2-H3.

\emph{Checks for deployers.} Re-setting \(\tau\) from genuine channel speech alone restores FRR but accepts 88--99.9\% of
held-out spoofs for the primary models; and at \(\tauE\) on ITW a transferred detector beats the better trivial system only
in narrow windows of bona fide priors (none for RawNet2).

% ============================ VI. RESULTS: ROUND 2 ======================================
\section{Round 2: Current Detectors, Retraining, Silence and Channels}
\label{sec:round2}
\subsection{Gates and operating behaviour}
\gentab{r2_gates}
\gentab{r2_operating}
\todo{Text from results/round2/analysis/r2\_explore (E1, E8, E10) and gates.json.}

\subsection{Registered tests}
\gentab{r2_tests}
\todo{R2-H1..H4 and the reading rule, from results/round2/analysis/r2\_registered.txt.}

\subsection{Seed variance and silence-trimmed training}
\todo{E2/E3 and R2-H2/H3.}

\subsection{Channels, languages and ASVspoof 5}
% src: results/round2/analysis/r2_registered.txt (R2-H4); numbers to be re-checked against the final run by the round-2 audit.
\emph{Channels (R2-H4).} H5 generalises to some channels and not to others. Low-bitrate speech codecs push genuine speech
far below the development threshold: Opus at 6\,kbps raises FRR by 40--69\,pp and AMR-NB at 4.75\,kbps by 77--92\,pp across
the four primary models, while FAR falls, so both meet R2-H4 for four of four models; Opus at 24\,kbps does too, with
shifts of at most 1\,pp in FRR. The phone impulse response raises FRR by 16--51\,pp (four of four). MP3 at 32\,kbps moves
almost nothing (one of four). Additive noise splits the models: AASIST and AASIST-L reject 41--74\,pp more genuine speech
at 10 and 5\,dB SNR, whereas RawNet2 and Res-TSSDNet accept more spoofs (Res-TSSDNet's FAR rises by 27--31\,pp), so the
noise conditions meet R2-H4 for two of four models only. R2-H4 holds for four of the eight conditions. \todo{E4 for the
round-2 detectors and the runs; E5 MLAAD-5; E6 ASVspoof~5.}

\subsection{Three-zone review and calibration across priors}
\todo{E7 and E9.}

% ============================ VII. DISCUSSION ==================================
\section{Discussion and Limitations}
\label{sec:discussion}
\todo{What carries over and what does not; recommendations for reporting (the card); triage; limitations (released
checkpoints vs recipes; one recipe retrained; ITW as the community's target; MLAAD's sample-rate confound; no adaptive
attackers; fp32 only).}

% ============================ VIII. CONCLUSION ==================================
\section{Conclusion}
\label{sec:conclusion}
\todo{Conclusion.}

\bibliographystyle{IEEEtran}
\bibliography{refs}

\end{document}
